\question
\begin{parts}
\part[5]
The following figure shows the range sizes for catarrhine primates in Africa.  The figure, representative for most taxa, illustrates a common biogeographic pattern. Briefly describe what this figure illustrates about the geographic range of most taxa.
\begin{figure}[h!]
\hfill \includegraphics{catarrhine.png}
\end{figure}
%\vspace{0.5in}

	\begin{solution}[0.5in]
	Small geographic range size.
	\end{solution}


\part[10]
Explain how the niche and competition with other species can influence the geographic range size of species.
	\begin{solution}
	Niche and competition reduce the  range size.
	\end{solution}

\end{parts}

\question[10]
Veter et al. (2013) examined whether Rapoport's rule applied to taxa that lived millions of years ago.  Below is a figure of their results.  Briefly explain Rapoport's Rule with regard to range size, then interpret their figure to explain whether their results are consistent with Rapoport's Rule. Explain each panel of the figure as it relates to the range size part of Rapoport's Rule.

\begin{center}
\includegraphics{paleorap.png} 
\end{center}

	\begin{AnswerPage}{2\basespace}
	Rapoport's rule says that geographic range size increase with higher latitudes. Also diversity decreases at higher latitudes (4 points).\vspace{\baselineskip}
	
	The Miocene and Pliocene show decreasing range size at higher latitudes so these periods are not consistent with Rapoport's Rule. The Pleistocene shows increasing range size at higher latitudes so that period is consistent with Rapoport. (2 points per panel).
	\end{AnswerPage}


\question[10]
What is the name for pattern shown in the four panels of the figure below. This pattern has been found for many taxa but does not appear to be universal. Explain two reasons why this pattern is not be universal for all taxa.



\question[5]
Explain how habitat quality can influence the range size of a species.

\AnswerBox{5\baselineskip}{%
Habitat quality\dots
}

\question[5] 
Explain how competition with other species can influence the geographic range size a species.

\AnswerBox{5\baselineskip}{%
Competition\dots
}

\question
Stevens (1996) tested whether Pacific coastal fishes followed Rapoport's Rule. His results (left panel) were consistent with the rule. Smith and Gaines (2003) analyzed Stevens' data in a different way. One part of their results is in the right panel.
\begin{parts}
\part[5]
Comparing the two panels, explain how Smith and Gaines analyzed the data differently to obtain different results, and tell whether their results supported Rappoport's Rule.

\bigskip
\begin{center}
\includegraphics{stevens1996} \hfill \includegraphics{smithgaines}
\end{center}

\bigskip

\AnswerBox{\basespace}{%
	Once they divided the data into different bathymetric depths, Rapoport's Rule no longer applied consistently.
}

\part[5]
What reason did Smith and Gaines give for the loss of species richness as shown in the right panel above?  Explain how this would cause the loss of richness.

\AnswerBox{\basespace}{%
	They proposed “loss of shallow water habitat type (coral reefs) from temperature latitudes.” The loss of coral reef habitat means you would lose the species richness associated with that habitat.
}

\end{parts}

